\subsection{弱空间的商空间与子空间}

设 F、G 为处于对偶的两个实向量空间。设 M 为 F 的一个向量子空间，并考察 G 中正交子空间 \(M^{\circ}\) 内的子空间 N；若 \(y_{1}, y_{2}\) 是 G 中模 N 同余的两点，则 \(\langle x, y_{1} \rangle = \langle x, y_{2} \rangle\) 对所有 \(x \in M\) 成立。对每个模 N 的类 \(\dot{y}\)，把 \(\langle x, y \rangle\) 当 y 取遍 \(\dot{y}\) 时所取的公共值记作 \(\langle x, \dot{y} \rangle\)；显然 \((x, \dot{y}) \mapsto \langle x, \dot{y} \rangle\) 是 \(M \times (G/N)\) 上的一个双线性型。

命题 7. --- 设 M 为 F 的一个向量子空间、N 为 G 的一个向量子空间，其中 F 与 G 是处于对偶的两个向量空间。假设 M 与 N 正交（这等价于 \(N \subset M^{\circ}\)，或 \(M \subset N^{\circ}\)）。则向量空间 M 与 G/N 按双线性型 \((x, \dot{y}) \mapsto \langle x, \dot{y} \rangle\) 处于对偶。

\begin{enumerate}
\def\labelenumi{(\roman{enumi})}
\item
  对这一对偶而言，拓扑 \(\sigma(\mathbf{M},\mathbf{G} / \mathbf{N})\) 由 \(\sigma(\mathbf{F},\mathbf{G})\) 诱导（特别地，我们有 \(\sigma(\mathbf{F},\mathbf{G}) = \sigma(\mathbf{F},\mathbf{G} / \mathbf{F}^{\circ})\)）。
\item
  对这一对偶而言，拓扑 \(\sigma(\mathbf{G}/\mathbf{N}, \mathbf{M})\) 粗于 \(\sigma(\mathbf{G}, \mathbf{F})\) 关于 \(\mathbf{N}\) 的商拓扑；这两个拓扑相同当且仅当 \(\mathbf{M} + \mathbf{G}^{\circ} = \mathbf{N}^{\circ}\)。
\setcounter{enumi}{0}
\item
  \(\mathbf{G} / \mathbf{N}\) 的每个元素都是模 \(\mathbf{N}\) 的一个类，由 \(\mathbf{G}\) 中的某个元素给出；若 \(z_{i}\)（\(1 \leqslant i \leqslant n\)）是 \(\mathbf{G}\) 的元素，且 \(\dot{z}_{i}\)（\(1 \leqslant i \leqslant n\)）是 \(z_{i}\) 在 \(\mathbf{G} / \mathbf{N}\) 中的类，则 \(y \in \mathbf{M}\) 中使 \(|\langle y, \dot{z}_{i} \rangle| \leqslant \alpha\) 对 \(1 \leqslant i \leqslant n\) 成立的 y 所成之集，是 \(\mathbf{M}\) 上这样的 \(x \in \mathbf{F}\)——即满足 \(|\langle x, z_{i} \rangle| \leqslant \alpha\)（\(1 \leqslant i \leqslant n\)）者——所成之集的迹；结论由弱拓扑零邻域的定义即得。
\item
  设 \(p: \mathbf{G} \to \mathbf{G} / \mathbf{N}\) 为典范满射。我们证明，商拓扑 \(\mathcal{T}\)——由 \(\sigma(\mathbf{G}, \mathbf{F})\) 按 \(\mathbf{N}\) 取商而得——与 \(\sigma(\mathbf{G} / \mathbf{N}, \mathbf{N}^{\circ})\) 一致。由于对 \(z \in \mathbf{G}\)、\(y \in \mathbf{N}^{\circ}\) 有 \(\langle y, p(z) \rangle = \langle y, z \rangle\)，可见 \(\sigma(\mathbf{G} / \mathbf{N}, \mathbf{N}^{\circ})\) 的每个零邻域都具有 \(p(\mathbf{V})\) 的形式，其中 \(\mathbf{V}\) 是 \(\sigma(\mathbf{G}, \mathbf{F})\) 的一个对关系 \(z - z' \in \mathbf{N}\) 饱和的零邻域，因此 \(\mathcal{T}\) 细于 \(\sigma(\mathbf{G} / \mathbf{N}, \mathbf{N}^{\circ})\)。 反之，设 \(\mathbf{U} = \mathbf{W}(y_1, ..., y_n; \alpha)\) 为 \(\mathbf{G}\) 中 \(\sigma(\mathbf{G}, \mathbf{F})\) 的一个零邻域，其中 \(y_i \in \mathbf{F}\)（\(1 \leqslant i \leqslant n\)），且 \(\alpha > 0\)；我们将看到，对 \(1 \leqslant i \leqslant n\)，存在元素 \(t_i \in \mathbf{N}^{\circ}\)，使得若令 \(\mathbf{U}' = \mathbf{W}(t_1, ..., t_n; \alpha)\)，则有 \(p(\mathbf{U}') \subset p(\mathbf{U})\)；这就表明 \(\sigma(\mathbf{G} / \mathbf{N}, \mathbf{N}^{\circ})\) 细于 \(\mathcal{T}\)，从而确实与 \(\mathcal{T}\) 一致。 现在，设 \(L\) 为 \(F\) 中由 \(N^{\circ}\) 与诸 \(y_i\) 生成的向量子空间，并用 \(P\) 表示 \(N^{\circ}\) 在 \(L\) 中的补子空间；它是有限维的，设其维数为 \(m\)。设 \((x_j)_{1 \leqslant j \leqslant m}\) 为 \(P\) 的一组基；限制到 \(N\) 上的线性型 \(x \mapsto \langle x_j, z \rangle\) 是线性无关的，因为否则便存在 \(x \neq 0\)（它属于 \(P\)），使得 \(\langle x, z \rangle = 0\) 对一切 \(z \in N\) 成立，也就是说 \(x \in N^{\circ}\)，这与 \(P\) 的定义矛盾。由此我们得出，对所有 \(z' \in G\)，存在 \(s \in N\) 使得 \(\langle x_j, z' \rangle = \langle x_j, s \rangle\) 对所有 \(j\) 成立；若 \(z' = z + s\)，则 \(\langle x, z \rangle = 0\) 对所有 \(x \in P\) 成立。作了这样的准备之后，写 \(y_i = t_i + w_i\)，其中 \(t_i \in N^{\circ}\)、\(w_i \in P\)；我们有 \(\langle y_i, z \rangle = \langle t_i, z \rangle = \langle t_i, z' \rangle\)（\(1 \leqslant i \leqslant n\)）；于是，对所有 \(z' \in U'\)，存在 \(z \in U\) 使得 \(z' - z \in N\)，也就是说有 \(p(\mathrm{U}') \subset p(\mathrm{U})\)。
\end{enumerate}

回到 M 是 \(N^{\circ}\) 中任意子空间的情形，注意显然有 \(\sigma(\mathrm{G}/\mathrm{N},\mathrm{M})=\sigma(\mathrm{G}/\mathrm{N},\mathrm{M}+\mathrm{G}^{\circ})\)；此外，由 II, p.~43 的命题 3 可见，若 \(y\in N^{\circ}\) 使得线性型 \(\dot{z}\mapsto\langle y,\dot{z}\rangle\) 关于 \(\sigma(\mathrm{G}/\mathrm{N},\mathrm{M})\) 连续，则必有 \(y\in M+G^{\circ}\)。我们由此断定：条件 \(M+G^{\circ}=N^{\circ}\) 是商拓扑 T 等于 \(\sigma(\mathrm{G}/\mathrm{N},\mathrm{M})\) 的充分必要条件。

注. --- M 与 G/N 之间的对偶（其中 M 与 N 是两个正交的子空间）在 M 中是分离的，当且仅当 \(M \cap G^{\circ} = \{0\}\)；它在 G/N 中是分离的，当且仅当 \(N = M^{\circ}\)。

推论 1. --- 假设 F 与 G 之间的对偶在 F 中是分离的。对 F 的向量子空间 M 而言，拓扑 \(\sigma(\mathbf{G}/\mathbf{M}^{\circ}, \mathbf{M})\) 与 \(\sigma(\mathbf{G}, \mathbf{F})\) 关于 \(\mathbf{M}^{\circ}\) 的商拓扑一致，当且仅当 M 关于拓扑 \(\sigma(\mathbf{F}, \mathbf{G})\) 是闭的。

这可由命题 7 取 \(\mathbf{N} = \mathbf{M}^{\circ}\)，并回忆 \(\mathbf{M}^{\circ \circ}\) 是 \(\mathbf{M}\) 关于 \(\sigma (\mathrm{F},\mathrm{G})\) 的闭包而得到（II, p.~45, 推论 2）。

推论 2. --- 若 M 是 n 维有限的且对偶在 F 中分离，则 \(M^{\circ}\) 在 G 中的余维数是 n。若 M 关于 \(\sigma(F, G)\) 是闭的且具有有限余维数 n，且对偶在 G 中分离，则 \(M^{\circ}\) 的维数是 n。

因为 \(G/M^{\circ}\) 与 M 处于分离对偶；若 M 的维数是 n，那么 \(G/M^{\circ}\) 亦然（II, p.~41, 例 1）。若 M 是闭的，则 \(F/M = F/M^{\circ\circ}\) 与 \(M^{\circ}\) 处于分离对偶；若 F/M 的维数是 n，那么 \(M^{\circ}\) 的维数也是 n（II, p.~41, 例 1）。

推论 3. --- 设 (F, G)、(F₁, G₁) 是两对处于分离对偶的空间，u 是 F 到 F₁ 中的线性映射，它关于 σ(F, G) 与 σ(F₁, G₁) 连续。则 u 是 F 到 F₁ 中的严格态射，当且仅当 Im(\({}^t u\)) 是 G 中关于 σ(G, F) 的闭子空间。

设 \(\mathbf{N} = \operatorname{Im}(^t u) \subset \mathbf{G}\)；我们知道 \(\mathrm{N}^{\circ} = \mathrm{Ker}(u)\) 在 \(\mathbf{F}\) 中成立（II, p.~47, 公式 (3)）。设 \(p: \mathbf{F} \to \mathbf{F}/\mathbf{N}^{\circ}\) 为典范映射，于是 \(u\) 分解为

\[
u: \mathrm{F} \xrightarrow {p} \mathrm{F} / \mathrm{N} ^ {\circ} \xrightarrow {w} \mathrm{F} _ {1},
\]

其中 w 是单射。空间 \(F/N^{\circ}\) 与 N 处于分离对偶，且由 II, p.~48 的公式 (1)，有 \(\langle w(\dot{y}), z_{1} \rangle = \langle \dot{y}, {}^{t}u(z_{1}) \rangle\)（对所有 \(\dot{y} \in F/N^{\circ}\) 与 \(z_{1} \in G_{1}\)）。这个关系式表明 w 是 \(F/N^{\circ}\)（赋拓扑 \(\sigma(F/N^{\circ}, N)\)）到 \(u(F)\)（赋 \(\sigma(F_{1}, G_{1})\) 诱导的拓扑）上的同构。因此结论由推论 1 以及严格态射的定义得出。

推论 4. --- 设 \((\mathrm{F}, \mathrm{G})\)、\((\mathrm{F}_{1}, \mathrm{G}_{1})\) 是两对处于分离对偶的空间，u 是 F 到 \(F_{1}\) 中的线性映射，它关于 \(\sigma(\mathrm{F}, \mathrm{G})\) 与 \(\sigma(\mathrm{F}_{1}, \mathrm{G}_{1})\) 连续。则 u 是满射，当且仅当 \(^{t}u\) 是 \(G_{1}\)（赋拓扑 \(\sigma(G_{1}, F_{1})\)）到 \(^{t}u(G_{1})\)（赋 \(\sigma(\mathrm{G}, \mathrm{F})\) 诱导的拓扑）上的同构。

因为 \(u(\mathrm{F}) = \mathrm{F}_{1}\) 等价于 \(u(\mathrm{F})\) 在 \(F_{1}\) 中关于 \(\sigma(\mathrm{F}_{1}, \mathrm{G}_{1})\) 既闭又处处稠密；于是推论 4 便由应用于 \(^{t}u\) 的推论 3 以及 II, p.~47 的推论 2 得出。

注. --- 1) 设 \((\mathrm{F}_1, \mathrm{G}_1), (\mathrm{F}_2, \mathrm{G}_2), (\mathrm{F}_3, \mathrm{G}_3)\) 为三对处于分离对偶的空间，考虑由两个线性映射组成的序列

\[
\mathrm{F} _ {1} \xrightarrow {u} \mathrm{F} _ {2} \xrightarrow {v} \mathrm{F} _ {3}\tag{5}
\]

它们分别关于相应于 \(\mathbf{G}_1\)、\(\mathbf{G}_2\)、\(\mathbf{G}_3\) 的弱拓扑连续；我们考虑转置映射的序列

\[
\mathrm{G} _ {3} \stackrel {{t _ {v}}} {{\to}} \mathrm{G} _ {2} \stackrel {{t _ {u}}} {{\to}} \mathrm{G} _ {1}.\tag{6}
\]

显然 \({}^t (v\circ u) = {}^t u\circ {}^tv\)，因此关系式 \(v\circ u = 0\) 等价于 \({}^t u\circ {}^tv = 0\)。序列 (5) 正合当且仅当下述三个条件成立：a) \({}^t u\circ {}^tv = 0\)；

\begin{enumerate}
\def\labelenumi{\alph{enumi})}
\setcounter{enumi}{1}
\item
  \(\operatorname{Im}(^{t}v)\) 在 \(\operatorname{Ker}(^{t}u)\) 中稠密；
\item
  \(^t u\) 是 \(G_2\) 到 \(G_1\) 中的严格态射。
\end{enumerate}

这事实上由 II, p.~49 的推论 3 以及 II, p.~47 的公式 (3) 与 (4) 得出。2) 切不可认为当 \(u\) 是 \(F\) 到 \(F_1\) 中的严格态射时，\(^t u\) 就必然是 \(G_1\) 到 \(G\) 中的严格态射；换言之，\(u\) 可以是严格态射，而 \(u(F)\) 却在 \(F_1\) 中关于 \(\sigma(F_1, G_1)\) 不闭。取 \(F\) 为 \(F_1\) 的一个非闭子空间且 \(G = G_1 / F^\circ\)、\(u\) 为典范单射，这一例子便说明了这一点。类似地，序列 (5) 正合这一事实并不必然蕴含 (6) 正合；然而，若序列 (5) 正合且 \(v\) 是严格态射，则序列 (6) 正合，这由注 1 以及 II, p.~49 的推论 3 得出。

\subsection{弱拓扑的乘积}

命题 8. --- 设 \((\mathbf{F}_{\iota}, \mathbf{G}_{\iota})_{\iota \in \mathbf{I}}\) 为一族处于对偶的空间对。设 \(\mathbf{F} = \prod_{\iota \in \mathbf{I}} \mathbf{F}_{\iota}\) 为诸 \(\mathbf{F}_{\iota}\) 的乘积空间，\(\mathbf{G} = \bigoplus_{\iota \in \mathbf{I}} \mathbf{G}_{\iota}\) 为诸 \(\mathbf{G}_{\iota}\) 的直和。若对所有 \(x = (x_{\iota}) \in \mathbf{F}\) 与所有 \(y = (y_{\iota}) \in \mathbf{G}\)，我们写 \(\langle x, y \rangle = \sum_{\iota \in \mathbf{I}} \langle x_{\iota}, y_{\iota} \rangle\)（这是一个只有有限多个非零项的和），则拓扑 \(\sigma(\mathbf{F}, \mathbf{G})\)（相对于双线性型 \((x, y) \mapsto \langle x, y \rangle\)）是诸拓扑 \(\sigma(\mathbf{F}_{\iota}, \mathbf{G}_{\iota})\) 的乘积。

因为，给定一个拓扑 \(\mathcal{T}\) 于 \(\mathbf{F}\) 上；为使对所有 \(y \in \mathbf{G}\) 线性型 \(x \mapsto \langle x, y \rangle\) 都关于 \(\mathcal{T}\) 连续，按 \(\langle x, y \rangle\) 的定义，其充分必要条件是：每个映射 \(x \mapsto \langle \mathrm{pr}_i x, y_i \rangle\) 都关于 \(\mathcal{T}\) 连续，其中 \(i\) 在 I 中任意，而 \(y_i\) 在 \(\mathbf{G}_i\) 中任意；而这意味着每个映射 \(\mathrm{pr}_i\)（F 到 \(\mathbf{F}_i\) 中）关于 \(\mathcal{T}\) 与 \(\sigma(\mathbf{F}_i, \mathbf{G}_i)\) 都连续（I, p.~10, 推论 1）；证毕。

注. --- F 与 G 之间的对偶在 F 中（相应地，在 G 中）分离，当且仅当对所有 \(\iota\in I\)，\(F_{\iota}\) 与 \(G_{\iota}\) 之间的对偶在 \(F_{\iota}\) 中（相应地，在 \(G_{\iota}\) 中）分离。若 F 与 G 之间的对偶在 F 中（相应地，在 G 中）分离，则在 F 中（相应地，在 G 中），与某个 \(G_{\iota}\)（相应地，\(F_{\iota}\)）正交的子空间——它被典范地等同于 G（相应地，F）的一个子空间——是诸 \(F_{\kappa}\)（\(\kappa\neq\iota\)）的乘积所成的子空间（相应地，诸 \(G_{\kappa}\)（\(\kappa\neq\iota\)）的直和）。

推论 1. --- 设 F 与 G 为两个处于分离对偶的向量空间。若空间 F（赋 \(\sigma(F, G)\)）是两个子空间 M、N 的拓扑直和，则空间 G（赋 \(\sigma(G, F)\)）是分别与 M 和 N 正交的子空间 M°、N° 的拓扑直和。

因为，设 \(p: \mathbf{F} \to \mathbf{M}\)、\(q: \mathbf{F} \to \mathbf{N}\) 为对应于 \(\mathbf{F}\) 分解为 \(\mathbf{M}\) 与 \(\mathbf{N}\) 之直和的投影；在这些条件下，映射 \((p, q): \mathbf{F} \to \mathbf{M} \times \mathbf{N}\) 是拓扑同构。若 \(\mathbf{M}_1 = \mathbf{G} / \mathbf{M}^\circ\)、\(\mathbf{N}_1 = \mathbf{G} / \mathbf{N}^\circ\)，则 \(\mathbf{M}\) 与 \(\mathbf{N}\) 上的拓扑（由 \(\mathbf{F}\) 的拓扑诱导）分别与 \(\sigma(\mathbf{M}, \mathbf{M}_1)\)、\(\sigma(\mathbf{N}, \mathbf{N}_1)\) 一致（II, p.~48, 命题 7）。映射 \(^t(p, q): \mathbf{M}_1 \times \mathbf{N}_1 \to \mathbf{G}\)，当给 \(\mathbf{M}_1, \mathbf{N}_1\) 与 \(\mathbf{G}\) 分别赋以拓扑 \(\sigma(\mathbf{M}_1, \mathbf{M})\)、\(\sigma(\mathbf{N}_1, \mathbf{N})\) 与 \(\sigma(\mathbf{G}, \mathbf{F})\) 时，由命题 8 是一个拓扑同构。在这个映射下，\(\mathbf{M}_1\)（相应地，\(\mathbf{N}_1\)）在 \(\mathbf{G}\) 中的像是子空间 \(\mathbf{N}^\circ\)（相应地，\(\mathbf{M}^\circ\)），而拓扑 \(\sigma(\mathbf{M}_1, \mathbf{M})\)（相应地，\(\sigma(\mathbf{N}_1, \mathbf{N})\)）的像是 \(\mathbf{N}^\circ\)（相应地，\(\mathbf{M}^\circ\)）上由 \(\sigma(\mathbf{G}, \mathbf{F})\) 诱导的拓扑，由此推论得证。

推论 2. --- 设 \((e_{\iota})_{\iota \in \mathbb{J}}\) 为向量空间 F 的一组基，F* 是 F 的对偶，并设 \(u: \mathbf{R}^{(I)} \to \mathbf{F}\) 为由这组基定义的（代数）同构。则转置映射 \({}^t u: \mathbf{F}^* \to \mathbf{R}^I\) 当 F* 赋拓扑 \(\sigma(\mathbf{F}^*, \mathbf{F})\)、\(\mathbf{R}^I\) 赋乘积拓扑时是拓扑同构。

我们知道（A, II, § 2.6, 命题 10），\({}^t u\) 是双射，并且若对某个 \(x^* \in \mathbf{F}^*\) 令 \(\langle e_i, x^* \rangle = \xi_i^*\)（对所有 \(i \in I\)），则像 \({}^t u(x^*)\) 是向量 \((\xi_i^*)\)，它属于 \(\mathbf{R}^I\)，从而对所有 \(x = \sum_{i} \xi_i e_i\)（在 \(F\) 中），有 \(\langle x, x^* \rangle = \sum_{i \in I} \xi_i \xi_i^*\)。推论随即由这个公式与命题 8 得出。

\subsection{弱完备空间}

命题 9. --- 设 F、G 为两个处于分离对偶的向量空间。若 \(\hat{\mathbf{F}}\) 是空间 F 关于拓扑 \(\sigma(\mathrm{F},\mathrm{G})\) 的完备化，且我们考虑典范单射 \(j:\mathrm{F}\to\mathrm{G}^{*}\)，其中 \(\mathrm{G}^{*}\) 赋拓扑 \(\sigma(\mathrm{G}^{*},\mathrm{G})\)，则连续延拓 \(\hat{j}:\hat{\mathrm{F}}\to\mathrm{G}^{*}\)（延拓自 \(j\)）是拓扑向量空间的一个同构。

因为，我们看到 \(G^{*}\) 赋以 \(\sigma(G^{*}, G)\) 后是 Hausdorff 的且完备（II, p.~51, 推论 2）；若借助 j 把 F 等同于 \(G^{*}\) 的一个向量子空间，则 \(\sigma(G^{*}, G)\) 在 F 上诱导的拓扑是 \(\sigma(F, G)\)，且 F 在 \(G^{*}\) 中关于拓扑 \(\sigma(G^{*}, G)\) 稠密（II, p.~43, 推论 4）；命题由此得证。

于是，关于某个弱拓扑完备的向量空间，就是任意的向量空间 G 的对偶 G* 赋以 \(\sigma(\mathbf{G}^{*}, \mathbf{G})\)；由 II, p.~51 的推论 2，它们（拓扑上）同构于实直线的乘积 \(R^{I}\)。为使措辞简化，我们将称它们为「直线积」（关于这类空间的内蕴刻画，见 II, p.~85, 习题 13 及 II, p.~81, 习题 1）。

我们注意到，在 \(G^{*}\) 上，拓扑 \(\sigma(G^{*}, G)\) 是 Hausdorff 的弱拓扑中最小的；因为，一个粗于 \(\sigma(G^{*}, G)\) 的弱拓扑必然形如 \(\sigma(G^{*}, H)\)，其中 \(H \subset G\)（II, p.~43, 推论 3）；但若 \(H \neq G\)，则存在一个线性型 \(x^{*} \in \mathbf{G}^{*}\)，它非零且与 \(\mathbf{H}\) 正交（A, II, § 7.3, 命题 8），因此 \(\sigma(\mathbf{G}^{*}, \mathbf{H})\) 不是 Hausdorff 的。

由这个注我们推出：若 F、G 是两个向量空间，则线性双射 \(u:G^{*}\to F^{*}\) 只要关于拓扑 \(\sigma(G^{*},G)\) 与 \(\sigma(F^{*},F)\) 连续，就必然是双连续的。

命题 10. --- 设 G 为实向量空间，F = G* 是它的对偶，赋拓扑 \(\sigma(\mathbf{G}^*, \mathbf{G})\)。

\begin{enumerate}
\def\labelenumi{(\roman{enumi})}
\item
  映射 \(\mathbf{V} \mapsto \mathbf{V}^{\circ}\) 是 \(\mathbf{G}\) 的向量子空间全体到 \(\mathbf{F}\) 的闭向量子空间全体上的双射。
\item
  \(\mathbf{F}\) 的每个闭向量子空间都是直线积且具有拓扑补。
\end{enumerate}

由双极定理（II, p.~45, 推论 2），\(\mathrm{V} \mapsto \mathrm{V}^{\circ}\) 是 G 中关于 \(\sigma(\mathbf{G}, \mathbf{G}^{*})\) 为闭的向量子空间 V 全体到 F 的闭向量子空间全体上的双射。但按定义，G 上每个线性型都关于 \(\sigma(\mathbf{G}, \mathbf{G}^{*})\) 连续，因此 G 中每个向量子空间都是闭的，因为它由一个方程组 \(\langle y, y_{\lambda}^{*} \rangle = 0\)（其中 \(y_{\lambda}^{*} \in \mathbf{G}^{*}\)）定义；这就证明了 (i)。

现在设 W 是 F 的一个闭子空间；于是有 \(W = V^{\circ}\)，其中 \(V = W^{\circ}\) 在 G 中。设 \(V'\) 是 V 在 G 中的一个补空间。我们知道 \(F = G^{*}\) 可典范地等同于 \(V^{*} \oplus V'^{*}\)，且 \(V'^{*}\) 等同于 \(V^{\circ} = W\)（A, II, § 2.6, 命题 10 的推论）；此外（II, p.~50, 命题 8），拓扑 \(\sigma(G^{*}, G)\) 可等同于拓扑 \(\sigma(V^{*}, V)\) 与 \(\sigma(V'^{*}, V')\) 的乘积；这就证明了断言 (ii)。

尽管对拓扑 \(\sigma(\mathbf{G},\mathbf{G}^{*})\) 而言 \(\mathbf{G}\) 的每个向量子空间都是闭的，我们仍注意到，若 \(\mathbf{G}\) 是无限维的，则拓扑 \(\sigma(\mathbf{G},\mathbf{G}^{*})\) 并不是 \(\mathbf{G}\) 上最细的局部凸拓扑，因为 \(\sigma(\mathbf{G},\mathbf{G}^{*})\) 的每个零邻域都包含一个无限维向量子空间：然而它是 \(\mathbf{G}\) 上弱拓扑中最细的（II, p.~43, 推论 3）。

\subsection{弱空间中的完备凸锥}

引理 1. --- 设 E 是一个 Hausdorff 弱空间，C 是 E 中以 0 为顶点的真锥，且它关于由 E 的一致结构所诱导的一致结构是完备的。则 E 上每个连续线性型都是两个在 E 上连续且在 C 上为正的线性型之差。

设 \(E'\) 为 E 的对偶，F 为 \(E'\) 的代数对偶，赋拓扑 \(\sigma(F, E')\)。设 \(H = C^{\circ} - C^{\circ}\) 为 \(E'\) 中由在 E 上连续且在 C 上为正的线性型之差所构成的向量子空间（II, p.~44, 命题 4）。只需证明 H 在 F 中的正交是 \(\{0\}\)（II, p.~41, 例 1）。于是设 \(a \in F\) 与 H 正交；由于 a 与 \(C^{\circ}\) 正交，它必属于 C 在 F 中的双极集。但 E 可等同于 F 的一个子空间，且由于 C 完备，从而在 F 中闭，故有 \(a \in C\)（II, p.~44, 定理 1）。类似地，a 与 \(-C^{\circ}\) 正交，因此 \(a \in -C\)。由于 C 是真锥，得 a = 0。

命题 11. --- 设 E 是一个 Hausdorff 弱空间，C 是 E 中以 0 为顶点的真凸锥，且它关于由 E 的一致结构所诱导的一致结构是完备的。

则存在一个集合 I 与一个 E 到乘积空间 \(R^{I}\) 中的连续线性映射 u，它具有下列性质：

\begin{enumerate}
\def\labelenumi{\alph{enumi})}
\item
  \(u\) 是 \(\mathbf{C}\) 到 \(u(\mathbf{C})\) 上的同构，两者分别赋以由 \(\mathbf{E}\) 的一致结构与 \(\mathbf{R}^{\mathrm{I}}\) 的一致结构诱导的一致结构。
\item
  我们有 \(u(\mathbf{C}) \subset \mathbf{R}_{+}^{\mathrm{I}}\)。
\end{enumerate}

此外，若 E 的一致结构在 C 上诱导的一致结构是可度量化的，则可取 I = N。

设 \((f_{\iota})_{\iota \in I}\) 为 E 上一族连续线性型，使得形如 \((x, y) \mapsto |f_{\iota}(x - y)|\) 的伪度量在 \(\mathbf{C} \times \mathbf{C}\) 上的有限和定义 C 的一致结构。（若该结构可度量化，则可取 \(\mathbf{I} = \mathbf{N}\)。）由引理 1，我们还可假设每个 \(f_{\iota}\) 在 C 上为正。设 \(u\) 为线性映射 \(x \mapsto (f_{\iota}(x))_{\iota \in I}\)，它把 E 映到 \(\mathbf{R}^{\mathrm{I}}\) 中。显然 \(u\) 连续且 \(u(\mathbf{C}) \subset \mathbf{R}_{+}^{\mathrm{I}}\)。限制 \(u|\mathbf{C}\) 是 C 到 \(u(\mathbf{C})\) 上的满的一致连续映射。此外，若 C 中的 \(x, y\) 使 \(f_{\iota}(x) = f_{\iota}(y)\) 对所有 \(\iota \in I\) 成立，则由于 C 的一致结构是 Hausdorff 的，有 \(x = y\)；因此 \(u|\mathbf{C}\) 是双射。最后，若 W 是 C 的一致结构的一个一致邻域，则存在 I 的有限子集 J 与数 \(\varepsilon > 0\)，使得关系式 \(|f_{\iota}(x) - f_{\iota}(y)| \leqslant \varepsilon\)（对 \(\iota \in J\)）蕴含 \((x, y) \in W\)；因此 \(u|\mathbf{C}\) 是 C 到 \(u(\mathbf{C})\) 上的、关于所考虑的一致结构的同构。

推论 1. --- 设 E 是一个 Hausdorff 弱空间，C 是 E 中以 0 为顶点的真凸锥，且它关于由 E 的一致结构所诱导的一致结构是完备的。则映射 \((x, y) \mapsto x + y\)（\(C \times C\) 到 C 中）是逆紧的。

由于命题 11，我们可以假设 \(\mathbf{E} = \mathbf{R}^{\mathbf{I}}\) 且 \(\mathbf{C} = \mathbf{R}_{+}^{\mathbf{I}}\)（GT, I, § 10.1, 推论 1 与 4）。但这时映射 \((x, y) \mapsto x + y\)（\(\mathbf{C} \times \mathbf{C}\) 到 \(\mathbf{C}\) 中）写作 \(((\xi_{\mathfrak{i}}), (\eta_{\mathfrak{i}})) \mapsto (\xi_{\mathfrak{i}} + \eta_{\mathfrak{i}})\)，于是我们只需证明连续映射 \(f: (\xi, \eta) \mapsto \xi + \eta\)（\(\mathbf{R}_{+} \times \mathbf{R}_{+}\) 到 \(\mathbf{R}_{+}\) 中）是逆紧的（GT, I, § 10.1, 推论 3）。现在，对所有 \(\zeta \in \mathbf{R}_{+}\)，我们看到 \(f(\zeta)\) 是诸对 \((\xi, \zeta - \xi)\)（满足 \(0 \leqslant \xi \leqslant \zeta\)）所成之集，因此 \(f\) 对区间 \([0, \zeta]\) 的逆像是 \((\xi, \eta) \in \mathbf{R}_{+} \times \mathbf{R}_{+}\) 中满足 \(\xi + \eta \leqslant \zeta\) 者所成之集，它是紧的。应用（GT, I, § 10.3, 命题 7）即得结论。

推论 2. --- 设 E 是一个 Hausdorff 弱空间，C 是 E 中以 0 为顶点的真凸锥，且它关于由 E 的一致结构所诱导的一致结构是完备的。

\begin{enumerate}
\def\labelenumi{(\roman{enumi})}
\item
  对每一点 \(a\)（属于 \(\mathbf{E}\)），交 \(\mathbf{C} \cap (a - \mathbf{C})\) 是紧的。
\item
  设 \(\mathbf{A},\mathbf{B}\) 是 C 中两个闭集。则 \(\mathbf{A} + \mathbf{B}\) 是 C 中的闭集。
\item
  由推论 1 及 GT, I, § 10.2, 定理 1, b)，\((x, y) \in \mathbf{C} \times \mathbf{C}\) 中满足 \(x + y = a\) 的那些所成之集是紧的。而这个集也就是诸 \((x, a - x)\)（其中 \(x \in \mathbf{C} \cap (a - \mathbf{C})\)）所成之集，这就证明了 (i)。
\item
  若 A 与 B 在 C 中闭，则 \(A \times B\) 在 \(C \times C\) 中闭，于是由推论 1 与 GT, I, § 10.1, 命题 1，\(A + B\) 在 C 中闭。
\end{enumerate}

\section{极点与极生成元}

\subsection{紧凸集的极点}

定义 1. --- 设 A 是仿射空间 E 中的一个凸集。若不存在既含于 A 又包含 x 的开线段，则称点 \(x \in A\) 为 A 的一个极点。

换言之，诸关系式 \(x=\lambda y+(1-\lambda)z\)、\(y\in A\)、\(z\in A\)、\(y\neq z\) 与 \(0\leqslant\lambda\leqslant1\) 蕴含 \(\lambda=0\) 或 \(\lambda=1\)（于是 x=y 或 x=z）。由此推出，x 不可能是 A 中带正质量的 n 个点 \(x_{i}\) 所成之集的重心，除非 x 是诸 \(x_{i}\) 之一；因为 n=2 时这正是定义；对任意的 n，则对 n 作归纳论证：由于 x 是 \(x_{1}\) 与 \(y_{1}\)——后者是诸 \(x_{i}\)（\(2\leqslant i\leqslant n\)）的重心——两者的重心，因此 x 与 \(x_{1}\) 或 \(y_{1}\) 重合，而在第二种情形只需应用归纳假设。

说 \(x\) 是 \(A\) 的极点，也等价于说 \(A - \{x\}\) 是凸的。

例. --- 1) 在空间 \(\mathbf{R}^n\) 中，球面 \(\mathbf{S}_{n-1}\) 的所有点都是闭球 \(\mathbf{B}_n\) 的极点。因为，若 \(\sum_{i} y_i^2 \leqslant 1\)、\(\sum_{i} z_i^2 \leqslant 1\) 且 \(0 < \lambda < 1\)，则关系式

\[
\lambda^ {2} \sum_ {i} y _ {i} ^ {2} + (1 - \lambda) ^ {2} \sum_ {i} z _ {i} ^ {2} + 2 \lambda (1 - \lambda) \sum_ {i} y _ {i} z _ {i} = 1 = (\lambda + (1 - \lambda)) ^ {2}
\]

仅在下述条件成立时才可能：

\[
\sum_ {i} y _ {i} ^ {2} = \sum_ {i} z _ {i} ^ {2} = \sum_ {i} y _ {i} z _ {i} = 1.
\]

但这蕴含 \(\sum_{i}(y_i - z_i)^2 = 0\)，从而 \(y_i = z_i\) 对所有 \(i\) 成立，这就证明了我们的断言。

\begin{enumerate}
\def\labelenumi{\arabic{enumi})}
\setcounter{enumi}{1}
\tightlist
\item
  在实数有界序列组成的赋范空间 \(\mathcal{B}(\mathbf{N})\)（I, p.~4）中，单位球的极点是点 \(x = (\xi_n)\)，它满足 \(|\xi_n| = 1\)（对所有 \(n\)）。因为，假设 \(|\xi_n| \leqslant 1\) 对所有 \(n\) 成立，而 \(|\xi_p| < 1\) 对某个指标 \(p\) 成立。这时我们可以写
\end{enumerate}

\[
x = \frac {1 + \xi_ {p}}{2} y + \frac {1 - \xi_ {p}}{2} z
\]

其中 \(y\)（相应地，\(z\)）是这样的点：它的所有坐标都等于 \(x\) 的同下标坐标，唯独下标 \(p\) 的坐标等于 1（相应地，- 1）。这表明 \(x\) 不是极点，因为我们有 \(\| y \| \leqslant 1\) 与 \(\| z \| \leqslant 1\)。反之，若 \(|\xi_n| = 1\) 对所有 \(n\) 成立，则 \(x\) 是极点，因为关系式 \(\xi_n = \lambda \eta_n + (1 - \lambda) \zeta_n\) 连同 \(|\eta_n| \leqslant 1\)、\(|\zeta_n| \leqslant 1\) 与 \(0 < \lambda < 1\) 蕴含 \(\xi_n = \eta_n = \zeta_n\)。

\begin{enumerate}
\def\labelenumi{\arabic{enumi})}
\setcounter{enumi}{2}
\tightlist
\item
  设 \(u: \mathrm{E} \to \mathrm{E}'\) 是仿射空间 \(\mathrm{E}\) 到仿射空间 \(\mathrm{E}'\) 中的仿射映射；设 \(\mathrm{C} \subset \mathrm{E}\)、\(\mathrm{C}' \subset \mathrm{E}'\) 是满足 \(u(\mathrm{C}) \subset \mathrm{C}'\) 的两个凸集。若 \(x'\) 是 \(\mathrm{C}'\) 的一个极点，且 \(x\) 是 \(u^{-1}(x') \cap \mathrm{C}\) 的一个极点，则 \(x\) 是 \(\mathrm{C}\) 的极点，这由定义 1 即得。
\end{enumerate}

命题 1. --- 设 B 是 A 的极点所成之集，其中 A 是 Hausdorff 局部凸空间 E 中的非空紧凸集，并设 f 是定义在 A 上且上半连续的凸函数。则 f 在 A 中的上确界可在 B 的（至少）一个点上达到。

用 F 表示由 A 的这样的子集 X 组成的族：X 非空、闭，且包含在 A 中并与 X 相交的每个开线段必含于 X。它具有下列性质；

\begin{enumerate}
\def\labelenumi{(\roman{enumi})}
\item
  A 属于 \(\mathfrak{F}\)。
\item
  点 \(a \in \mathbf{A}\) 使 \(\{a\} \in \mathfrak{F}\)，当且仅当 \(a\) 是 \(\mathbf{A}\) 的一个极点。
\item
  每个非空交 \(\mathbf{X}\)（族 \((\mathbf{X}_{\alpha})\) 的交，该族由 \(\mathfrak{F}\) 中的集组成）也属于 \(\mathfrak{F}\)。
\end{enumerate}

性质 (i)、(ii) 与 (iii) 由定义立得。

\begin{enumerate}
\def\labelenumi{(\roman{enumi})}
\setcounter{enumi}{3}
\tightlist
\item
  设 \(X \in F\)，h 是 A 中凸且上半连续的函数；则 X 中使限制 \(h|X\) 在 X 中达到其上确界的点所成的集 Y 满足：Y 属于 F。
\end{enumerate}

因为，\(h|X\) 在 \(X\) 上上半连续，故其上确界 \(\alpha\)（在 \(X\) 上取的上确界）在 \(X\) 的至少一个点上达到（GT, IV, § 6.2, 定理 3）；于是 \(Y\) 非空，它也是闭的（GT, IV, § 6.2, 命题 1）。另一方面，设 \(x, y\) 是 \(A\) 的两个不同点，并设 \(z = \lambda x + (1 - \lambda)y\) 是 \(Y\) 中满足 \(0 < \lambda < 1\) 的一点；由于 \(Y \subset X\) 且 \(X \in \mathfrak{F}\)，我们有 \(x \in X\) 与 \(y \in X\)；另一方面，由于 \(h\) 是凸的，我们有

\[
h (z) \leqslant \lambda h (x) + (1 - \lambda) h (y)
\]

但由于 \(h(x) \leqslant \alpha\)、\(h(y) \leqslant \alpha\) 且 \(h(z) = \alpha\)，必有 \(h(x) = h(y) = \alpha\)，也就是说 \(x \in \mathbf{Y}\) 且 \(y \in \mathbf{Y}\)。因此 \(\mathbf{Y} \in \mathfrak{F}\)。

在建立了这些性质之后，设 M 是使 f 在 A 中达到其上确界的 \(x \in A\) 所成之集；由 (iv)，\(M \in F\)。另一方面，由 (iii) 以及 F 中的集都是紧集 A 的闭子集这一事实，可知 F 关于序关系 \(\supset\) 是归纳的。由 S, III, §2.4 的定理 2，M 包含一个子集 N，它是 F 的极小元。我们将证明 N 由单独一点组成，从而完成命题的证明。由于 E 是 Hausdorff 局部凸空间，只需证明 E 上每个连续线性型 u 在 N 中是常数（II, p.~38, 推论 1）。而由 (iv) 推出，集 \(N'\)——由 \(x \in N\) 中使 \(u|N\) 在 N 中达到其上确界者组成——满足 \(N'\) 属于 F；由于 N 在 F 中极小，必有 \(N' = N\)。

推论. --- 设 A 是 Hausdorff 局部凸空间 E 中的紧凸集。则 A 的每个闭支撑超平面 H 至少包含 A 的一个极点。

因为，若 \(f(x) = \alpha\) 是 \(H\) 的一个方程，且 \(f(x) \leqslant \alpha\) 在 \(A\) 中成立，只需对 \(f\) 应用命题 1。

定理 1（Krein-Milman）. --- 在 Hausdorff 局部凸空间 E 中，每个紧凸集 A 都是它的极点所成之集的闭凸包。

因为，设 C 是 A 的极点所成之集的闭凸包；显然 \(C \subset A\)。为看出 \(A \subset C\)，只需证明：若 u 是在 E 中连续的仿射线性函数，且在 C 中 \(u(x) \geqslant 0\)，则在 A 中也有 \(u(x) \geqslant 0\)（II, p.~39, 推论 4）；而这由对 -u 应用命题 1 即得。

命题 2. --- 设 x 是 Hausdorff 局部凸空间 E 中紧凸集 A 的一个极点。则对 E 中 x 的每个开邻域 V，存在 E 中的一个开半空间 F，使得 \(x \in F \cap A \subset V \cap A\)（换言之，包含 x 的诸开半空间在 A 上的迹构成 x 在 A 中的基本邻域系）。

对 E 的每个包含 x 的开半空间 D，集 \(A \cap \overline{D}\) 是 x 在 A 中的一个紧邻域，而所有这些邻域的交恰为点 x（任意两个不同的点可被一个闭超平面严格分离（II, p.~38, 命题 4）。由 GT, I, § 9.2 的命题 1，只需证明诸集 \(A \cap \overline{D}\) 构成一个滤子基。现在，若写 \(L_{D} = A \cap (E - D)\)，则集 \(L_{D}\) 是凸的、紧的，且含于凸集 \(A - \{x\}\)；若 \(D_{1}, D_{2}\) 是 E 的两个包含 x 的开半空间，则 \(L_{D_{1}} \cup L_{D_{2}}\) 的凸包 B 因此含于 \(A - \{x\}\)；但 B 是紧集（II, p.~14, 命题 15），故存在一个闭超平面 H 把 x 与 B 严格分离（II, p.~38, 命题 4），而若由 H 决定且包含 x 的开半空间是 D，则我们有 \(L_{D_{1}} \cup L_{D_{2}} \subset L_{D}\)，因此 \(A \cap \overline{D} \subset (A \cap \overline{D}_{1}) \cap (A \cap \overline{D}_{2})\)。

推论. --- 在 Hausdorff 局部凸空间中，设 K 是紧凸集 A 的一个紧子集。则下列条件等价。

\begin{enumerate}
\def\labelenumi{\alph{enumi})}
\item
  A 是 K 的闭凸包。
\item
  K 与每一集相交，该集是 A 与其某一支撑超平面的交。
\item
  K 包含 A 的极点所成之集。
\end{enumerate}

\(a) \Rightarrow b)\)。假设存在 A 的一个支撑超平面 H，其方程为 \(f(x) = \alpha\)，使得 \((H \cap A) \cap K = \varnothing\)，并例如假设在 A 中 \(f(x) \geqslant \alpha\)。由于按假设 \(f(x) - \alpha > 0\) 对所有 \(x \in K\) 成立，且 K 是紧的，我们有

\[
\beta = \inf _ {x \in \mathbf {K}} f (x) > \alpha ,
\]

于是 K 含于闭半空间 \(f(x) \geqslant \beta\)；因此对 K 的闭凸包 A 同样如此，而这是荒谬的。

\(b) \Rightarrow c)\)。假设 \(x\) 是 \(\mathbf{A}\) 的一个不属于 \(\mathbf{K}\) 的极点；则存在 \(\mathbf{V}\)（\(x\) 在 \(\mathbf{E}\) 中的一个邻域），使得 \(\mathbf{V} \cap \mathbf{A} \cap \mathbf{K} = \varnothing\)。但由命题 2，我们可以假设 \(\mathbf{V}\) 是由超平面 \(\mathbf{H}\)（其方程为 \(f(z) = \alpha\)）所定义的开半空间。若例如 \(f(x) > \alpha\)，则对所有 \(y \in \mathbf{K}\)，我们有 \(f(y) \leqslant \alpha\)，因此 \(\mathbf{K}\) 不与 \(\mathbf{A}\) 和支撑超平面 \(f(z) = \gamma > \alpha\)（平行于 \(\mathbf{H}\)）的交相遇（II, p.~37, 命题 2）；这是荒谬的。

\(c) \Rightarrow a)\)。这是 Krein-Milman 定理的明显推论。

习题 11）。定理 1（II, p.~56）的证明表明，无论如何，A 都是 A 中属于某一支撑超平面的那些极点所成之集的闭凸包。

\subsection{凸锥的极生成元}

设 C 是向量空间 E 中以 0 为顶点的一个凸锥；显然 C 中除顶点外没有别的点能是极点；顶点是 C 的极点，当且仅当 C 是真的且尖的。

定义 2. --- 设 C 是向量空间 E 中以 0 为顶点的一个凸锥。若每条含于 C、不含 0 且与 D 相交的开线段都含于 D，则称一条从 0 出发的半线 D ⊂ C 为 C 的一个极生成元。

这等价于说：对所有 \(x \in D\)（满足 \(x \neq 0\)），若 \(y \neq 0\)、\(y' \neq 0\) 是 C 中满足 \(x = y + y'\) 的两个点，则必有 \(y \in D\) 与 \(y' \in D\)。

注 1. --- 设 C 是 E 中的一个真尖凸锥，并在 E 上考虑使 C 成为 \(\geqslant 0\) 元素集的序结构（II, p.~12, 命题 13）；为使 E 的一个元素——设为 x \textgreater{} 0——属于 C 的某个极生成元，其必要充分条件是：每个以 x 为上界的元素 \(y \geqslant 0\) 都具有形式 \(\lambda x\)，其中 \(0 \leqslant \lambda \leqslant 1\)：事实上，说 y 以 x 为上界，意味着 \(x = y + y'\)，其中 \(y' \in C\)，由此即得结论。

命题 3. --- 在向量空间 E 中，设 C 是以 0 为顶点的凸锥，\(x_{0} \neq 0\) 是 C 的一点，D 是含于 C、从 0 出发且包含 \(x_{0}\) 的一条半线。设 H 是包含 \(x_{0}\) 且不通过 0 的一个超平面。则 D 是 C 的极生成元，当且仅当 \(x_{0}\) 是 H ∩ C 的一个极点。

该条件显然是必要的。反之，设它成立；假设存在一条直线 \(D'\)，它不含 D、通过 \(x_{0}\)，且使 \(D' \cap C\) 包含一个含 \(x_{0}\) 的开线段。设 \(y \neq 0\) 是 \(D'\) 的一个方向向量；诸假设蕴含：点 \((1 + \lambda)x_{0} + \mu y\) 当 \(|\lambda|\) 与 \(|\mu|\) 充分小时属于 C。但在由 D 与 \(D'\) 所决定且赋有典范拓扑的平面 P 中，\(x_{0}\) 是 \(P \cap C\) 的一个内点，由此推出，直线 \(P \cap H\) 包含一个含于 \(H \cap C\) 且含 \(x_{0}\) 的开线段。这与假设矛盾。

定义 3. --- 设 C 是 Hausdorff 拓扑向量空间 E 中的一个凸集。C 的一个非空紧凸集 A 称为 C 的一个帽，如果 A 在 C 中的补集 C−A 是凸的。

设 C 是 E 中以 0 为顶点的一个尖凸锥，并设 A 是 C 的一个帽。写 B = C - A。对每条从 0 出发的闭半线 L ⊂ C，集 L ∩ A 与 L ∩ B 都是凸集，它们在 L 中互为补集，其并为 L，且 L ∩ A 是紧的。由于 L ∩ A 至少对一条半线 L 非空，我们看出 0 ∈ A，从而 L ∩ A 是以 0 为一个端点的闭线段。若 A 存在，则 C 是真锥。

命题 4. --- 设 C 是 E（一个 Hausdorff 局部凸空间）中以 0 为顶点的一个尖凸锥。

\begin{enumerate}
\def\labelenumi{\alph{enumi})}
\tightlist
\item
  设 A 是 C 的一个帽。设 p 是 A 的度规到 C 上的限制（II, p.~20）。\(x \in C\) 中使 \(p(x) \leqslant 1\) 者所成之集就是集 A。函数 p 是下半连续的，且具有下列性质：
\end{enumerate}

\begin{enumerate}
\def\labelenumi{(\roman{enumi})}
\item
  对任意的 \(x, y\)（在 \(\mathbf{C}\) 中），我们有 \(p(x + y) = p(x) + p(y)\)。
\item
  对任意的 \(x \in \mathbf{C}\) 与 \(\lambda \in \mathbf{R}_{+}^{*}\)，我们有 \(p(\lambda x) = \lambda p(x)\)。
\item
  若 \(x \in C\)，则关系 \(p(x) = 0\) 等价于 x = 0。
\end{enumerate}

\begin{enumerate}
\def\labelenumi{\alph{enumi})}
\setcounter{enumi}{1}
\tightlist
\item
  反之，设 \(p\) 是定义在 \(\mathbf{C}\) 上、取值于 \([0, +\infty]\) 的函数，满足 a) 的条件 (i)、(ii)。设 \(\mathbf{A}\) 为由 \(x \in \mathbf{C}\) 中使 \(p(x) \leqslant 1\) 者组成的集。则 \(\mathbf{A}\) 与 \(\mathbf{C} - \mathbf{A}\) 都是凸的。A 是帽，当且仅当 \(\mathbf{A}\) 紧且非空。
\end{enumerate}

陈述 \(b)\) 是明显的。\(a)\) 中所述的性质是命题 4 前面的那些注记以及 II, p.~20 的命题 22 与 II, p.~20 的命题 23 的推论，但下式除外：

\[
p (x + y) \geqslant p (x) + p (y).
\]

只需对 \(x \neq 0\) 且 \(y \neq 0\) 的情形证明最后这个不等式；于是有 \(p(x) > 0\)、\(p(y) > 0\)。设 \(\mu, \lambda\) 是两个 \textgreater{} 0 的数，满足 \(\lambda < p(x)\)、\(\mu < p(y)\)，并用 B 表示 A 在 C 中的补集。我们有 \(x \in \lambda B\)、\(y \in \mu B\)，因此 \(x + y \in \lambda B + \mu B\)；由 B 的凸性，有 \(\lambda B + \mu B \subset (\lambda + \mu)B\)，从而 \(p(x + y) > \lambda + \mu\)，这就蕴含上面所述的不等式。

推论 1. --- 设 C 是 E（一个 Hausdorff 局部凸空间）中以 0 为顶点的尖凸锥，并设 p 是 C 的一个帽 A 的度规。则 A 的极点就是点 0，以及 C 的诸极生成元上满足 \(p(x) = 1\) 的那些点 x。

显然 0 是 A 的极点。设 x 是 C 的一个极生成元 L 上的点，且满足 \(p(x) = 1\)。设 y、z 是 A 中满足 \(x = \frac{1}{2}(y + z)\) 的两点。由于 L 是极的，我们有 \(y = \lambda x\) 与 \(z = \mu x\)，其中 \(\lambda\) 与 \(\mu\) 是 \(\geqslant 0\) 的数，满足 \(\frac{1}{2}(\lambda + \mu) = 1\)、\(\lambda = \lambda p(x) = p(y) \leqslant 1\) 与 \(\mu = \mu p(x) = p(z) \leqslant 1\)，由此得 \(\lambda = \mu = 1\)，从而 y = z = x；所以 x 是 A 的极点。反之，设 \(x \neq 0\) 是 A 的一个极点。显然 \(p(x) = 1\)。设 y、\(y'\) 是 C 中满足 \(x = y + y'\) 的两点，我们将证明 y、\(y'\) 与 x 成比例。不失一般性，可以假设数 \(\lambda = p(y)\) 与 \(\lambda' = p(y')\) 有限且 \textgreater0。于是由命题 4 (i)，\(\lambda^{-1}y \in A\)、\(\lambda'^{-1}y' \in A\)、\(\lambda + \lambda' = 1\)，且等式 \(x = \lambda(\lambda^{-1}y) + \lambda'(\lambda'^{-1}y')\) 按假设蕴含

\[
x = \lambda^ {- 1} y = \lambda^ {\prime - 1} y ^ {\prime}.
\]

推论 2. --- C 中属于 C 的某个帽的每一点，也属于 C 的诸极生成元之并的闭凸包。

这由推论 1 与 Krein-Milman 定理（II, p.~55, 定理 1）立得。

* 例. --- 设 X 是一个局部紧且 \(\sigma\) -紧的空间。设 C 是 \(\mathcal{M}_{+}(X)\)（赋淡拓扑）中以 0 为顶点的一个闭凸锥。我们将证明 C 是它的诸帽之并。设 \((X_{n})\) 是 X 的开相对紧集的一个递增序列，其并为 X。设 \(\mu\) 是 C 中 \(\neq 0\) 的一个元素。存在 \(\alpha_{n} > 0\) 使 \(\sum \alpha_{n}\mu(X_{n}) = 1\)。

对每个测度 \(v \in C\)，令 \(p(v) = \sum_{n} \alpha_{n} v(X_{n}) \in [0, +\infty]\)。C 上的函数 p 满足命题 4 的条件 (i) 与 (ii)。它对淡拓扑是下半连续的（INT, IV, 第 2 版, § 1, No.~1, 命题 4）。于是 \(\gamma \in C\) 中使 \(p(\gamma) \leqslant 1\) 者所成的集 A 是闭的且非空。另一方面，X 的每个紧集都含于某个 \(X_{n}\) 之中，于是 A 既为淡有界，从而也是淡紧的（INT, III, 第 2 版, § 1, No.~9, 命题 15）。因此集 A 是 C 的一个包含 \(\mu\) 的帽。

命题 5. --- 设 C 是 E（一个 Hausdorff 弱空间）中以 0 为顶点的一个真凸锥；假设 C 对由 E 的一致结构所诱导的一致结构是完备的，并且 C 中 0 有一个可数的基本邻域系。则 C 是它的诸帽之并，并且是其诸极生成元之并的闭凸包。

第二个断言由第一个断言及上面的推论 2 得出。利用 II, p.~52 的命题 11，可把命题化归为 \(\mathbf{E} = \mathbf{R}^{\mathrm{I}}\) 且 \(\mathbf{C} \subset \mathbf{R}_{+}^{\mathrm{I}}\) 的情形。对所有 \(\alpha \in \mathbf{I}\)，把投影 \(\operatorname{pr}_{\alpha}\)（\(\mathbf{E}\) 中的）记作 \(f_{\alpha}\)；于是 \(f_{\alpha}\) 是连续线性型。另一方面，设 \((\mathrm{V}_n)_{n \in \mathbb{N}}\) 是 \(\mathbb{C}\) 中 0 的一个可数基本邻域系。按 \(\mathbf{E}\) 的拓扑的定义，对每个 \(n \in \mathbb{N}\)，存在一个有限子集 \(\mathrm{J}_n\)（\(\mathbf{I}\) 的）与一个数 \(\varepsilon_n > 0\)，使得 \(\mathrm{V}_n\) 包含 \(\mathrm{W}_n\)——即 \(x \in \mathbb{C}\) 中使 \(f_{\alpha}(x) \leqslant \varepsilon_n\)（对所有 \(\alpha \in \mathrm{J}_n\)）者所成的集；令 \(\mathrm{J} = \bigcup_{n \in \mathbb{N}} \mathrm{J}_n\)。设 \(y \neq 0\) 是 \(\mathbb{C}\) 中的一点，并设 \(p\) 是函数 \(\sum_{\alpha \in \mathbf{J}} \lambda_{\alpha}(f_{\alpha}|\mathbb{C})\)，其中诸 \(\lambda_{\alpha} > 0\) 取得使 \(p(y) = 1\)；这是可能的，因为若 \(f_{\alpha}(y) = 0\) 对所有 \(\alpha \in \mathbf{J}\) 成立，则 \(y \in \mathrm{V}_n\) 对所有 \(n\) 成立，这就蕴含 \(y = 0\)，与假设相违。现在我们注意到，对所有 \(\alpha \in \mathbf{I}\)，函数 \(f_{\alpha}|\mathbb{C}\) 在点 0 处连续，因此存在一个 \(n \in \mathbb{N}\)，使 \(f_{\alpha}\) 在某个 \(\mathrm{W}_n\) 中有界，从而在 \(\mathbb{C}\) 中以有限多个函数 \(f_{\beta}|\mathbb{C}\)（其中 \(\beta \in \mathbf{J}\)）的一个线性组合为其一个上界。由此推出，若 \(\mathbf{A}\) 是由 \(x \in \mathbb{C}\) 中使 \(p(x) \leqslant 1\) 者组成的集，则 \(f_{\alpha}\) 在 \(\mathbf{A}\) 中有界（对所有 \(\alpha \in \mathbf{I}\)）。由于 \(p\) 在 \(\mathbb{C}\) 中下半连续，可知 \(\mathbf{A}\) 在 \(\mathbb{C}\) 中闭且非空，从而是紧的。既然显然 \(p\) 满足 II, p.~58 命题 4 的条件 (i) 与 (ii)，我们就看出 \(\mathbf{A}\) 是 \(\mathbb{C}\) 中的一个帽，且包含 \(y\)。

注 2. --- 存在弱完备且没有极生成元的真凸锥（II, p.~92, 习题 31）。

\subsection{具紧底的凸锥}

命题 6. --- 设 E 是一个 Hausdorff 局部凸空间，K 是 E 中不包含 0 的一个紧凸集。则以 0 为顶点且包含 K 的最小尖锥 C 是 E 中的一个真凸锥，并且是 E 的一个局部紧且完备的子空间；此外，存在 E 中一个不包含 0 的闭超平面 H，它与含于 C 的所有从 0 出发的半线都相交，且使 H ∩ C 是紧的。进而，若 D 是由具有这些性质的一个闭超平面 H 所决定的包含 0 的半空间，则 C ∩ D 是 C 的一个帽，且 C 是诸 λ(C ∩ D)（λ \textgreater{} 0）之并。

由 II, p.~38 的命题 4，存在一个把 0 与 K 严格分离的闭超平面 H。现在，\(\{0\}\) 与 K 之并的凸包 A 是紧的（II, p.~14, 命题 15），并且是诸 \(\lambda K\)（\(0 \leqslant \lambda \leqslant 1\)）之并。由于 0 与 K 严格地位于 H 的两侧，对每个 \(x \in K\)，存在 \(\lambda\) 满足 \(0 < \lambda < 1\) 且 \(\lambda x \in H\)。由于 C 是诸 \(\lambda A\)（\(\lambda \geqslant 1\)）之并，我们看出 H 与含于 C 的每条从 0 出发的半线都相交，且 \(H \cap A = H \cap C\) 是紧的。进而，C 也是诸 \(\lambda(H \cap C)\)（\(\lambda \geqslant 0\)）之并；设 \(C_n\) 为诸 \(\lambda(H \cap C)\)（\(0 \leqslant \lambda \leqslant n\)）之并。显然 \(C_n\) 是 \(\{0\}\) 与 \(n(H \cap C)\) 之并的凸包，因此是紧的。而且，对所有 \(x \in E\)，存在 x 在 E 中的一个闭邻域 V 和一个整数 n，使 \(V \cap C \subset C_n\)；事实上，若 H 由方程 \(f(z) = \alpha\)（其中 \(\alpha > 0\)）定义，只需对 V 取由 nH 决定且包含 0 的闭半空间，其中 n 大到使 \(n\alpha > f(x)\)。这表明 C 是局部紧的（取 \(x \in C\)），并且它在 E 中是闭的。我们还可以把 K 看作完备化 \(\hat{E}\) 的一个子集，因此 C 在 \(\hat{E}\) 中也是闭的，从而是完备的。

给定 Hausdorff 拓扑向量空间 E 中的一个锥 C 和一个闭超平面 H，使得 H 不含 C 的顶点 s，且 C 是以 s 为顶点、包含 \(H \cap C\) 的最小锥，则我们称交 \(H \cap C\) 为锥 C 的一个「底」。命题 6 表明，在 Hausdorff 局部凸空间 E 中，以 0 为顶点且包含一个不含 0 的紧凸集 K 的最小锥是具有紧底的锥，并且每个具有紧底 S 的凸锥都是局部紧且完备的。

例. --- 1) 有限维向量空间 E 中的每个真的闭凸锥都具有紧底。事实上，由 II, p.~52 命题 11，只需考虑 \(E = R^{n}\) 且 \(C = R_{+}^{n}\) 的情形。若 \((e_{i})_{1 \leqslant i \leqslant n}\) 是 \(R^{n}\) 的典范基，则显然以诸 \(e_{i}\)（\(1 \leqslant i \leqslant n\)）的凸包为紧凸集者是 \(R_{+}^{n}\) 的一个紧底。

* 2) 若 X 是紧空间，则 X 上正测度所成的锥 \(\mathcal{M}_{+}(\mathrm{X})\)（赋淡拓扑）是具有紧底的锥（INT, III, 第 2 版, § 1, No.~9, 命题 15 的推论 3）。*

\section{复局部凸空间}

\subsection{\texorpdfstring{\(\mathbf{C}\) 上的拓扑向量空间}{复数域上的拓扑向量空间}}

设 E 是复数域 C 上的一个拓扑向量空间；E 的拓扑也与把标量域限制到 R 所得的 R 上向量空间结构相容。我们用 \(E_{0}\) 表示 E 底下的 R 上拓扑向量空间（I, p.~2）。注意，在 \(E_{0}\) 中，映射 \(x \mapsto ix\)（它不是一个位似）是 \(E_{0}\) 的拓扑向量结构的一个自同构 u，满足 \(u^{2}(x) = -x\)。

反之，设 \(\mathbf{F}\) 是 \(\mathbf{R}\) 上的一个拓扑向量空间，并假设存在一个自同构 \(u\)（\(\mathbf{F}\) 上的），使 \(u^2 = -1_{\mathbf{F}}(1_{\mathbf{F}}\) 是 \(\mathbf{F}\) 的恒等自同构。我们知道（A, IX, § 3, No.~2），这时可以在 \(\mathbf{F}\) 上定义一个关于 \(\mathbf{C}\) 的向量空间结构：写 \(\lambda x = \alpha x + \beta u(x)\)，它对所有 \(\lambda = \alpha + i\beta \in \mathbf{C}\) 与所有 \(x \in \mathbf{F}\) 成立。此外，由于映射 \((\alpha, \beta, x) \mapsto \alpha x + \beta u(x)\)（从 \(\mathbf{R}^2 \times \mathbf{F}\) 到 \(\mathbf{F}\) 中）是连续的，\(\mathbf{F}\) 的拓扑与上面定义的关于 \(\mathbf{C}\) 的向量空间结构相容；若用 \(\mathbf{E}\) 表示这样定义的 \(\mathbf{C}\) 上的拓扑向量空间，则 \(\mathbf{F}\) 就是 \(\mathbf{R}\) 上那个作为 \(\mathbf{E}\) 之底的拓扑向量空间。

注. --- 给定 R 上的一个拓扑向量空间 F，并非总有平方为 \(-1_{F}\) 的自同构 u 存在；例如，在有限奇数维的 R 上向量空间上，不可能定义关于 C 的向量空间结构。

设 E 是 C 上的一个拓扑向量空间，\(E_{0}\) 是 E 底下的 R 上拓扑向量空间。E 中的每个线性流形 M 也是 \(E_{0}\) 中的线性流形，但反之不成立。为避免混淆，我们称关于 C（相应地，关于 R）的向量空间结构的线性流形为复（相应地，实）线性流形。有限维 n（相应地，有限余维 n）的复线性流形是维数 2n（相应地，余维数 2n）的实线性流形。为使 E 的实向量子空间 M 同时也是复向量子空间，必要且充分的是 \(iM \subset M\)。

回忆一下，若 E 与 F 是 C 上的两个拓扑向量空间，则称 E 到 F 中的一个映射为 C-线性的（相应地，R-线性的），如果它对 E 与 F 的关于 C（相应地，R）的向量空间结构是线性映射；每个 C-线性映射显然是 R-线性的，但反之不成立。我们称 E 上的 C-线性型为复线性型，称 E 上的 R-线性型（即 \(E_{0}\) 上的线性型）为实线性型。若 f 是 E 上的一个复线性型，则显然 f 的实部 g = Rf 与虚部 h = If 都是实线性型；此外，关系 \(f(ix) = if(x)\) 蕴含恒等式 \(h(x) = -g(ix)\)；换言之，我们有

\[
f (x) = (\mathscr {R} f) (x) - i (\mathscr {R} f) (i x).\tag{1}
\]

反之，若 g 是 E 上的一个实线性型，则 \(f(x) = g(x) - ig(ix)\) 是 E 上使 Rf = g 的唯一的复线性型；而 f 连续当且仅当 g 连续。

现在设 H 是 E 中的一个复超平面，其方程为 \(f(x) = \alpha + i\beta\)，其中 f 是 E 上的一个复线性型；令 \(g = Rf\)，我们看出 H 是两个实超平面 \(H_{1}\)、\(H_{2}\) 的交，它们的方程分别为 \(g(x) = \alpha\) 与 \(g(ix) = -\beta\)；若 H 是闭的，则 \(H_{1}\) 与 \(H_{2}\) 也是闭的（I, p.~13, 定理 1）。反之，设 \(H_{0}\) 是一个齐次实超平面，其方程为 \(g(x) = 0\)（其中 g 是 E 上的一个实线性型）；则 H——即 \(H_{0}\) 与 \(iH_{0}\) 的交——是一个齐次复超平面，而若 f 是使 Rf = g 的那个复线性型，则 \(f(x) = 0\) 便是 H 的方程；若 \(H_{0}\) 是闭的，则 H 也是闭的。

设 G 是 R 上的一个拓扑向量空间，并设 \(\mathbf{G}_{(\mathbf{C})}\) 是把标量域扩张到 C 后由 G 得到的 C 上向量空间（A, II, § 5.1）。把 G 等同于 \(\mathbf{G}_{(\mathbf{C})}\) 的一个子集，借助映射 \(x \mapsto 1 \otimes x\)。于是，R-线性映射 \((x, y) \mapsto x + i \cdot y\) 是 \(G \times G\) 到 \(\mathbf{G}_{(\mathbf{C})}\) 上的一个双射，借助它，我们把 \(G \times G\) 的乘积拓扑搬到 \(\mathbf{G}_{(\mathbf{C})}\) 上。这样，赋有此拓扑的 \(\mathbf{G}_{(\mathbf{C})}\) 是 C 上的一个拓扑向量空间。我们称 \(\mathbf{G}_{(\mathbf{C})}\) 为 G 的复化的拓扑向量空间。

\subsection{复局部凸空间}

说复向量空间 E 的子集 A 是均衡的，意思是：对所有 \(x \in A\)，有 \(\rho x \in A\)（对 \(0 \leqslant \rho \leqslant 1\)），且 \(e^{i\vartheta}x \in A\) 对所有实数 \(\vartheta\) 成立。

我们说 E 的一个集合 A 是凸的，如果它在 E 底下的实空间 \(E_{0}\) 中是凸的。为使 E 的凸集 \(A \neq \varnothing\) 是均衡的，只需 \(e^{i\vartheta}A \subset A\) 对所有实数 \(\vartheta\) 成立；因为这首先蕴含 -A = A；由 A 是凸的，我们看出 0 属于 A，从而 \(\rho A \subset A\)（对 \(0 \leqslant \rho \leqslant 1\)）。

设 E 是一个复拓扑向量空间。包含 E 的集合 A 的最小均衡凸集（相应地，闭均衡凸集）称为 A 的均衡凸包（相应地，闭均衡凸包）；A 的闭均衡凸包是 A 的均衡凸包的闭包。后者是诸集 \(e^{i\theta}A\) 之并的凸包；因此我们可以把它定义为诸线性和 \(\sum_{i}\lambda_{i}x_{i}\) 所成的集，其中 \((x_{i})\) 是 A 的点的任意有限族，而 \((\lambda_{i})\) 是满足 \(\sum_{i}|\lambda_{i}| \leqslant 1\) 的复数族。若 A 是预紧的，则其均衡包也是预紧的（I, p.~6, 命题 3）。

我们称复拓扑向量空间 E 是局部凸的，如果 E 底下的实拓扑向量空间 \(E_{0}\) 是局部凸的，也就是说，如果 E 中 0 的每个邻域都包含 0 的一个凸邻域；E 上的一个拓扑 T 称为局部凸的，如果它与 E 的（关于 C 的）向量空间结构相容，并且赋有拓扑 T 的 E 是局部凸的。由于在这种情形，0 的每个闭凸邻域 V 都包含一个均衡邻域 W（I, p.~7, 命题 4），我们看出 V 还包含 U，即 W 的闭均衡凸包；换言之，0 的均衡闭凸邻域构成 E 中 0 的基本邻域系，它在每个比率 \(\neq 0\) 的位似下不变。

反之，设 E 是一个复向量空间，并设 S 是 E 上由吸收的均衡凸集组成的一个滤子基。这时我们知道（II, p.~23, 命题 1），由 S 中的集经比率为 \textgreater{} 0 的位似变换得到的诸集所成之集 B，构成 E 底下的实向量空间 \(E_{0}\) 上某个局部凸拓扑 T 的 0 的基本邻域系。此外，由于 B 中的集都是均衡的，它们在每个位似 \(x \mapsto e^{i\vartheta}x\) 下不变，这表明 T 与 E 的（C 上的）向量空间结构相容（I, p.~7, 命题 4）。

复向量空间 E 上的每个局部凸拓扑都可以由一组半范数定义，因为 0 的一个开均衡凸邻域的度规是 E 上的一个半范数。

对实局部凸空间详细给出的、载于 II, p.~25 至 36 的概念与结果，都可毫无修改地推广到复局部凸空间，只需把对称凸集换成均衡凸集。

复局部凸空间若可度量化且完备，就称为 Fréchet 空间。

\subsection{Hahn-Banach 定理及其应用}

定理 1（Hahn-Banach）. --- 设 V 是复向量空间 E 的一个向量子空间，f 是 V 上的一个（复）线性型，p 是 E 上的一个半范数，满足 \(|f(y)| \leqslant p(y)\)（对所有 \(y \in V\)）。则存在 E 上的一个线性型 \(f_{1}\)，它延拓 f 且满足 \(|f_{1}(x)| \leqslant p(x)\)（对所有 \(x \in E\)）。

因为 \(g = \mathcal{R}f\) 是定义在 \(V\) 中的实线性型，且满足 \(|g(y)| \leqslant p(y)\) 于 \(V\) 的每一点；因此存在一个实线性型 \(g_1\) 于 \(E\) 中，它延拓 \(g\) 且满足 \(|g_1(x)| \leqslant p(x)\)（对所有 \(x \in E\)）（II, p.~23, 推论 1）。设 \(f_1(x) = g_1(x) - ig_1(ix)\) 为 \(E\) 上以 \(g_1\) 为实部的那个复线性型（II, p.~61）。对所有实数 \(\vartheta\)

\[
\left| \mathcal {R} (e ^ {i \vartheta} f _ {1} (x)) \right| = \left| \mathcal {R} (f _ {1} (e ^ {i \vartheta} x)) \right| = \left| g _ {1} (e ^ {i \vartheta} x) \right| \leqslant p (e ^ {i \vartheta} x) = p (x)
\]

由于 \(p\) 是复空间 \(\mathbf{E}\) 上的一个半范数；这蕴含关系 \(|f_1(x)| \leqslant p(x)\)，定理得证。

推论 1. --- 设 \(x_0\) 是复拓扑向量空间 \(E\) 中的一点，\(p\) 是 \(E\) 中的一个连续半范数；则存在一个连续（复）线性型 \(f\)，它定义在 \(E\) 中，使得 \(f(x_0) = p(x_0)\)，且 \(|f(x)| \leqslant p(x)\)（对所有 \(x \in E\)）。

推论 2. --- 设 V 是复局部凸空间 E 的一个向量子空间，f 是 V 中定义且连续的一个（复）线性型；则存在 E 中定义且延拓 f 的一个连续线性型 \(f_{1}\)。若 E 是赋范的，则存在这样的线性型 \(f_{1}\)，它还满足 \(\|f_{1}\|=\|f\|\)。

推论 3. --- 设 M 是 Hausdorff 复局部凸空间 E 的一个有限维向量子空间。则存在 E 的一个闭向量子空间 N，它是 M 在 E 中的拓扑余子空间。

这些利用 p.~24 定理 1 的证明，与 II, p.~23 推论 2、p.~24 推论 3、p.~24 命题 2 以及 p.~25 推论 2 中的证明相同。

命题 1. --- 设 A 是复拓扑向量空间 E 中的一个非空开凸集，M 是一个不与 A 相交的非空（复）线性流形。则存在一个包含 M 且不与 A 相交的闭复超平面 H。

我们可以假设 \(0 \in M\)。于是存在一个包含 M 且不与 A 相交的闭实超平面 \(H_{0}\)（II, p.~36; 定理 1）。由于 M = iM，闭复超平面 \(\mathbf{H} = \mathbf{H}_{0} \cap (i\mathbf{H}_{0})\) 就有所要求的性质。

推论. --- 在复局部凸空间 E 中，每个闭复线性流形 M 都是包含它的那些闭复超平面的交。

事实上，对所有 \(x \notin M\)，存在 x 的一个不与 M 相交的开凸邻域 V，从而存在一个包含 M 且不与 V 相交的闭复超平面 H；更强的，H 不包含 x。

命题 2. --- 设 A 是复拓扑向量空间 E 的一个非空均衡开凸集，B 是一个不与 A 相交的非空凸集。则存在 E 上的一个连续复线性型 f 和一个数 \(\alpha > 0\)，使得 \(|f(x)| < \alpha\) 在 A 中成立，且 \(|f(y)| \geqslant \alpha\) 在 B 中成立。

因为，存在连续实线性型 \(g\) 于 \(\mathbf{E}\) 上，以及一个实数 \(\alpha\)，使得 \(g(x) < \alpha\) 在 \(\mathbf{A}\) 中成立，且 \(g(y) \geqslant \alpha\) 在 \(\mathbf{B}\) 中成立（II, p.~37, 命题 1）。由于 \(0 \in \mathbf{A}\)，我们有 \(\alpha > 0\)。我们证明，连续复线性型 \(f(x) = g(x) - ig(ix)\) 与数 \(\alpha\) 具有所要求的性质。因为，由于 \(\mathcal{R}f = g\)，我们有 \(|f(y)| \geqslant \alpha\) 在 \(\mathbf{B}\) 中成立。另一方面，对所有 \(x \in \mathbf{A}\) 与所有实数 \(\vartheta\)，点 \(e^{i\vartheta}x\) 属于 \(\mathbf{A}\)（因为 \(\mathbf{A}\) 是均衡的），且我们有 \(f(x) = e^{-i\vartheta}f(e^{i\vartheta}x)\)；于是存在一个数 \(\vartheta\)，使得 \(|f(x)| = \mathcal{R}(e^{i\vartheta}f(x)) = g(e^{i\vartheta}x) < \alpha\)，命题得证。

命题 3. --- 设 A 是复局部凸空间 E 中的一个均衡闭凸集，K 是 E 中一个不与 A 相交的非空紧凸集。则存在 E 上的一个连续复线性型 f 和一个数 \(\alpha > 0\)，使得 \(|f(x)| < \alpha\) 在 A 中成立，且 \(|f(y)| > \alpha\) 在 K 中成立。

本命题由 II, p.~38 命题 4 得出，正如命题 2 由 II, p.~37 命题 1 得出一样。

\subsection{复向量空间上的弱拓扑}

II, § 6 No. 1 与 No. 2 的定义和结果可不加改变地适用于复向量空间。若 F 与 G 是由双线性型 B 处于对偶的两个复向量空间，则底空间 \(F_{0}\) 与 \(G_{0}\) 由 RB 处于对偶，且由 II, p.~61 公式 (1) 推出，弱拓扑 \(\sigma(\mathrm{F},\mathrm{G})\) 与 \(\sigma(\mathrm{F}_{0},\mathrm{G}_{0})\) 是相同的。

定义 1. --- 设 F 与 G 是处于对偶的两个复向量空间。对 F 的任意子集 M，M 在 G 中的极集，记作 \(M^{\circ}\)，是 \(y \in G\) 中使 \(\mathcal{R}(\langle x, y \rangle) \geqslant -1\)（对所有 \(x \in M\)）者所成的集。

若 \(\mathbf{M}^{\circ}\) 是 \(\mathbf{M} \subset \mathbf{F}\) 在 \(\mathbf{G}\) 中的极集，则 \((\lambda \mathbf{M})^{\circ} = \lambda^{-1}\mathbf{M}^{\circ}\) 对所有 \(\lambda \in \mathbf{C}^{*}\) 成立。

若 M 是 E 的一个（复）向量子空间，则 \(M^{\circ}\) 是一个闭向量子空间（对 \(\sigma(\mathbf{G},\mathbf{F})\) 而言），因为关系 \(\mathcal{R}(\lambda\langle x,y\rangle)\geqslant-1\)（对每个标量 \(\lambda\in C\)）蕴含 \(\langle x,y\rangle=0\)；我们仍说 \(M^{\circ}\) 是 G 中与 M 正交的子空间。

若 M 是 F 中的一个均衡集，则 \(M^{\circ}\) 是 G 中的一个均衡集；这时 \(M^{\circ}\) 是 \(y \in G\) 中使 \(\left|\langle x, y \rangle\right| \leqslant 1\)（对所有 \(x \in M\)）者所成的集；因为这个关系等价于 \(\mathcal{R}(\langle \zeta x, y \rangle) \leqslant 1\) 对所有 \(x \in M\) 与所有 \(\zeta \in C\)（满足 \(|\zeta| = 1\)）成立。

II, p.~41 至 51 的结果对复向量空间也无限制地成立。

\exercisesheading
